\chapter{曲线}\label{ch:b2:curves}

几何开始转向\omterm{def:b2:diffcalc:differential}{微分}。一条曲线是空间中运动的点，
微积分给出它的速度与加速度；几何关心的是内蕴的量，
即不依赖于沿轨迹运动快慢的量。答案之一是\emph{弧长}，
它度量轨迹本身；另一个是\emph{\omterm{thm:b2:curves:frenet2d}{曲率}}，它度量曲线如何弯曲。
在 $3$ 维情形下，另一个不变量\emph{\omterm{thm:b2:curves:frenet3d}{挠率}}，
度量曲线如何扭出所在平面。
描述这一切的
基本工具是随之运动的\emph{弗雷内标架}。

\section{参数弧}

\begin{definition}[参数弧]\label{def:b2:curves:arc}
称\emph{参数弧}\index{参数弧}为
$\mathcal{C}^k$（$k \geq 1$）类的
映射 $\gamma \colon I \to \R^n$，它为
$\mathcal{C}^k$ 类，定义于区间 $I$ 上。点 $\gamma(t)$ 称为\emph{正则点}，若
$\gamma'(t) \neq 0$；弧上各点皆为正则点时，称该弧为\emph{正则弧}。
过 $\gamma(t)$、以 $\gamma'(t)$ 为方向的直线
称为正则点处的\emph{切线}\index{切线}。
\end{definition}

\begin{definition}[参数变换]\label{def:b2:curves:reparam}
称\emph{参数变换}\index{参数变换}为
$\mathcal{C}^k$ 类的
$\mathcal{C}^k$ \omterm{def:b2:diffcalc:differential}{微分}同胚 $\theta \colon J \to I$（两区间之间），
且处处 $\theta' \neq 0$。称 $\gamma$ 与
$\gamma \circ \theta$ 为\emph{等价}的；一条\emph{几何弧}\index{几何弧}
（或称曲线）即一个等价类。
在参数变换下不变的概念——
如轨迹、\omterm{def:b2:curves:arc}{切线}、
弧长、\omterm{thm:b2:curves:frenet2d}{曲率}——称为\emph{几何的}。
\end{definition}

\begin{remark}
仅凭轨迹并不能确定\omterm{def:b2:curves:reparam}{几何弧}：
参数化 $t \mapsto (\cos t, \sin t)$ 在 $[0, 2\pi]$ 与
$[0, 4\pi]$ 上的像相同，却分别绕圆一周与两周。
\omterm{def:b2:curves:reparam}{几何弧}记住的是遍历的重数与方向，
而非速度。
\end{remark}

\begin{example}[速度变化不影响
几何性质]\label{ex:b2:curves:speedinvariance}
将单位圆参数化为
\[
\gamma(t) = (\cos t^2,\ \sin t^2),
\qquad t \in \intcc0{\sqrt{2\pi}} .
\]
速率 $\norm{\gamma'(t)} = 2t$ 线性增长，而
\[
L = \int_0^{\sqrt{2\pi}}2t\,\dd t = 2\pi ,
\]
与匀速时的长度相同，
正如\cref{thm:b2:curves:lengthinv}所言，借助\omterm{def:b2:curves:reparam}{参数变换} $t \mapsto t^2$。
\omterm{def:b2:curves:arc}{切线}、
由\cref{prop:b2:curves:kappaformula}算得的\omterm{thm:b2:curves:frenet2d}{曲率}，以及其他一切
几何量也都一致；只有 $t = 0$
值得留意，此处 $\gamma'(0) = 0$ 使\emph{此参数化}
成为非正则的，尽管其轨迹恰为一个圆。
几何陈述很难容忍糟糕的参数化：
先重新参数化，
再作结论。
\end{example}

\begin{example}\label{ex:b2:curves:cusp}
弧 $\gamma(t) = (t^2, t^3)$ 为 $\mathcal{C}^\infty$ 类，但不是
正则的：$\gamma'(0) = (0, 0)$。其轨迹，即半立方抛物线
$y^2 = x^3$，在原点处有一个\emph{尖点}：参数化的光滑性
并不能阻止几何奇点的出现，
即速度为零之处。这正说明正则性假设
$\gamma' \neq 0$ 并非形式上的要求。
\end{example}

\section{弧长}

\begin{definition}[弧长]\label{def:b2:curves:length}
设 $\gamma \colon [a, b] \to \R^n$ 为 $\mathcal{C}^1$ 弧。其
\emph{弧长}\index{弧长}为
\[
L(\gamma) = \int_a^b \norm{\gamma'(t)}\, \dd t ,
\]
其中 $\norm{\cdot}$ 为欧几里得\omterm{def:b2:nvs:norm}{范数}。称\emph{弧长函数}\index{弧长函数}
为以 $t_0$ 为基点的 $s(t) = \int_{t_0}^t
\norm{\gamma'(u)}\,\dd u$。
\end{definition}

\begin{theorem}[弧长是几何量；折线
刻画]\label{thm:b2:curves:lengthinv}
\leavevmode
\begin{enumerate}
\item 若 $\theta \colon [c, d] \to [a, b]$ 为 $\mathcal{C}^1$
\omterm{def:b2:curves:reparam}{参数变换}，则 $L(\gamma \circ \theta) = L(\gamma)$。
\item $L(\gamma)$ 是内接折线长度的
上确界：
\[
L(\gamma)
= \sup\Bigl\{\, \sum_{i=1}^{m} \norm{\gamma(t_i) -
\gamma(t_{i-1})} \;:\; a = t_0 < t_1 < \dots < t_m = b \,\Bigr\}.
\]
\end{enumerate}
\end{theorem}

\begin{proof}
\emph{1.} 由变量替换 $t = \theta(u)$（参见本科一年级分册，
因 $\theta$ 为 $\mathcal{C}^1$ 单调映射，该替换有效）：
\[
\int_c^d \norm{(\gamma\circ\theta)'(u)}\,\dd u
= \int_c^d \norm{\gamma'(\theta(u))}\,\abs{\theta'(u)}\,\dd u
= \int_a^b \norm{\gamma'(t)}\,\dd t ,
\]
其中用到了 $(\gamma\circ\theta)' = \theta'\cdot
(\gamma'\circ\theta)$，并且若 $\theta$ 递减，
则 $\theta'$ 的符号可由积分上下限的交换吸收。

\emph{2.} 对任意分割，$\gamma(t_i) - \gamma(t_{i-1}) =
\int_{t_{i-1}}^{t_i}\gamma'(t)\,\dd t$，由积分形式的三角不等式
可得 $\norm{\gamma(t_i) - \gamma(t_{i-1})}
\leq \int_{t_{i-1}}^{t_i}\norm{\gamma'}$：任一折线的长度都不超过
$L(\gamma)$，于是 $\sup \leq L(\gamma)$。

为证反向不等式，设 $\varepsilon > 0$。由于 $\gamma'$
在紧集 $[a,b]$ 上\omterm{def:b2:metric:continuity}{连续}，故它一致\omterm{def:b2:metric:continuity}{连续}：
存在 $\delta > 0$，使得 $\norm{\gamma'(t) - \gamma'(u)} \leq
\varepsilon$，只要 $\abs{t - u} \leq \delta$。取一个
步长 $\leq \delta$ 的分割。对每个小区间，当 $t \in [t_{i-1}, t_i]$ 时，
\[
\gamma(t_i) - \gamma(t_{i-1})
= \int_{t_{i-1}}^{t_i}\gamma'(t)\,\dd t
= (t_i - t_{i-1})\,\gamma'(t_{i-1}) + R_i,
\qquad
\norm{R_i} \leq \varepsilon\,(t_i - t_{i-1}),
\]
由于 $R_i = \int_{t_{i-1}}^{t_i}(\gamma'(t) -
\gamma'(t_{i-1}))\,\dd t$。由此得
\[
\norm{\gamma(t_i) - \gamma(t_{i-1})}
\geq (t_i - t_{i-1})\norm{\gamma'(t_{i-1})}
- \varepsilon (t_i - t_{i-1}) .
\]
求和，并将 $\sum (t_i -
t_{i-1})\norm{\gamma'(t_{i-1})}$ 与 $\int_a^b\norm{\gamma'}$ 比较（它是
\omterm{def:b2:metric:continuity}{连续}函数 $\norm{\gamma'}$ 的黎曼和，
与积分之差在 $\varepsilon(b - a)$ 之内，只需 $\delta$ 充分小，
这再次由一致\omterm{def:b2:metric:continuity}{连续}性保证），于是得到一条长度 $\geq
L(\gamma) - 2\varepsilon(b - a)$ 的折线。令 $\varepsilon \to 0$
即证得结论。
\end{proof}

\begin{example}[阿基米德与内接
多边形]\label{ex:b2:curves:archimedes}
对单位圆，内接正 $n$ 边形的长度为
$L_n = 2n\sin\frac\pi n$，而展开式 $\sin x = x -
\frac{x^3}6 + O(x^5)$ 给出
\[
L_n = 2\pi - \frac{\pi^3}{3n^2} + O\Bigl(\frac1{n^4}\Bigr) :
\]
\cref{thm:b2:curves:lengthinv} 中的折线长度
以\emph{二次}速度收敛。数值上，$L_6 = 6$
（即正六边形，给出粗略的 $\pi > 3$），而 $L_{96} =
192\sin\frac{\pi}{96} \approx 6.28206$，与 $2\pi \approx
6.28319$——误差 $0.00113$ 与预测值
$\pi^3/(3\cdot96^2) \approx 0.00112$。这正是阿基米德能够
将正六边形\omterm{def:b2:metric:continuity}{连续}加倍五次至 $96$ 边，仅凭手算便能把
$\pi$ 夹到三位数字精度：每加倍一次，
误差便缩为原来的四分之一。多边形刻画中的上确界
不只是在极限处达到，而是
\emph{迅速}地达到，因为光滑曲线与其弦
只在二阶项上才分开。
\end{example}

\begin{theorem}[弧长参数化]\label{thm:b2:curves:arclength}
设 $\gamma \colon I \to \R^n$ 为正则 $\mathcal{C}^k$ 弧
（$k \geq 1$）。\omterm{def:b2:curves:length}{弧长函数} $s$ 是 $\mathcal{C}^k$
\omterm{def:b2:diffcalc:differential}{微分}同胚，它把 $I$ 映到区间 $J$ 上，且 $\tilde\gamma =
\gamma \circ s^{-1}$ 满足 $\norm{\tilde\gamma'} = 1$
处处成立。在相差参数平移与定向的意义下，
这种\emph{弧长}（或\emph{单位速率}）参数化是
唯一的。
\end{theorem}

\begin{proof}
$s'(t) = \norm{\gamma'(t)} > 0$，且 $s'$ 是 $\mathcal{C}^{k-1}$
（$\mathcal{C}^{k-1}$ 映射 $\gamma'$ 与\omterm{def:b2:nvs:norm}{范数}的复合，
在 $0$ 之外光滑），故 $s$ 是 $\mathcal{C}^k$，严格
递增，为到 $J = s(I)$ 上的双射，且其逆
为 $\mathcal{C}^k$ 的，这由维数为 $1$ 时的反函数定理
（本科一年级分册）。于是
\[
\tilde\gamma'(\sigma)
= \frac{\gamma'(t)}{s'(t)}
= \frac{\gamma'(t)}{\norm{\gamma'(t)}},
\qquad t = s^{-1}(\sigma),
\]
为单位向量。若 $\hat\gamma = \gamma\circ\theta$ 为另一
单位速率参数化，则 $\abs{\theta'} = 1$，故 $\theta'
= \pm 1$ 为常数（由\omterm{def:b2:metric:continuity}{连续}性），即 $\theta(u) = \pm u + c$。
\end{proof}

\begin{remark}
弧长是区分几何与
动力学的参数。一条轨迹可以配以任意的
速率分布走过——（即运动的物理）——，但每一个
与参数化无关的问题（形状、弯曲、
密切性）都有一个标准的“时钟”，即所走过的距离。
这就是为什么下文所有\omterm{thm:b2:curves:frenet2d}{曲率}公式都先
在单位速率下\emph{定义}，然后再借助链式法则
\emph{转换}到任意参数化：转换因子
是 $v = s'$ 的幂，正确追踪它们正是
\cref{prop:b2:curves:kappaformula} 的全部内容。
\end{remark}

\begin{example}[圆与螺旋线]\label{ex:b2:curves:helix}
对圆 $\gamma(t) = (R\cos t, R\sin t)$，$\norm{\gamma'} =
R$，故 $s = Rt$，一整圈的长度为 $2\pi R$。对
螺旋线 $\gamma(t) = (a\cos t,\ a\sin t,\ bt)$，其中 $a > 0$，
$\norm{\gamma'(t)} = \sqrt{a^2 + b^2}$ 为常数：该螺旋线
以恒定速率走过，且 $s = t\sqrt{a^2 + b^2}$。
\end{example}

\begin{example}[极坐标中的
弧长]\label{ex:b2:curves:polarlength}
极坐标曲线 $r = r(\theta)$ 即弧 $\gamma(\theta) =
(r\cos\theta,\ r\sin\theta)$，其中
\[
\gamma'(\theta) = (r'\cos\theta - r\sin\theta,\
r'\sin\theta + r\cos\theta),
\qquad
\norm{\gamma'(\theta)}^2 = r'^2 + r^2
\]
（交叉项相消）：极坐标下的弧长元素为
$\sqrt{r^2 + r'^2}\,\dd\theta$。对心形线 $r = 1 +
\cos\theta$:
\[
r^2 + r'^2 = (1 + \cos\theta)^2 + \sin^2\theta = 2 +
2\cos\theta = 4\cos^2\tfrac\theta2,
\]
而在 $\intcc{-\pi}{\pi}$ 上，半角余弦
非负，故
\[
L = \int_{-\pi}^{\pi}2\cos\tfrac\theta2\,\dd\theta
= \Bigl[4\sin\tfrac\theta2\Bigr]_{-\pi}^{\pi}
= 4 - (-4) = 8 :
\]
与 \cref{exo:b2:curves:1} 中的摆线拱相似，由圆
构成的曲线长度为有理数，式中不含 $\pi$
，处处如此。半角分解是计算由圆\omterm{def:b2:structures:generated}{生成}曲线长度的标准技巧；
当它失效时（如
椭圆），长度是一个真正的新函数——一个
椭圆积分，超出了初等闭式的范围。
\end{example}

\section{平面曲率}

本节中，弧均为 $\mathcal{C}^2$ 的正则弧，
位于有向欧氏平面中。以弧长为参数，
记 $T(s) = \tilde\gamma'(s)$ 为单位切向量，并记 $N(s)$
为与 $T(s)$ 正交的单位向量（$T$ 经
$+\pi/2$ 旋转所得）。

\begin{theorem}[平面弗雷内公式]\label{thm:b2:curves:frenet2d}
设 $\tilde\gamma$ 为有向平面中的单位速率 $\mathcal{C}^2$ 弧，
则存在\omterm{def:b2:metric:continuity}{连续}函数 $\kappa$，
称其为（代数）\emph{曲率}\index{曲率}，使得
\[
T'(s) = \kappa(s)\, N(s),
\qquad
N'(s) = -\kappa(s)\, T(s).
\]
\end{theorem}

\begin{proof}
由于 $\norm{T(s)}^2 = 1$ 对所有 $s$ 成立，对
数量积求导得 $2\langle T'(s), T(s)\rangle = 0$，故 $T'(s)$ 与
$T(s)$ 正交，从而与 $N(s)$ 共线（维数为 $2$）；记 $T'(s)
= \kappa(s) N(s)$，其中 $\kappa(s) = \langle T'(s), N(s)\rangle$，
\omterm{def:b2:metric:continuity}{连续}。同理 $N' \perp N$，故 $N' = \lambda T$。
再对 $\langle T, N\rangle = 0$ 求导，得 $\langle T', N\rangle +
\langle T, N'\rangle = \kappa + \lambda = 0$。
\end{proof}

\begin{definition}\label{def:b2:curves:curvature}
当 $\kappa(s) \neq 0$ 时，
称\emph{曲率半径}\index{曲率半径}为 $R(s) =
1/\abs{\kappa(s)}$，而
\emph{曲率中心}\index{曲率中心}为
$\tilde\gamma(s) + \frac{1}{\kappa(s)} N(s)$；以此为圆心、
$R(s)$ 为半径的圆称为\emph{密切圆}\index{密切圆}，
即曲线在 $\tilde\gamma(s)$ 处的最佳圆近似。
\end{definition}

\begin{example}[指数函数的
密切圆]\label{ex:b2:curves:oscexp}
对 $y = \eu^x$ 在点 $(0, 1)$ 处：$f'(0) = f''(0) = 1$，
故由下面的图像公式，
\[
\kappa(0) = \frac{1}{(1 + 1)^{3/2}} = \frac1{2\sqrt2},
\qquad R = 2\sqrt2 .
\]
单位切向量为 $T = \frac{(1, 1)}{\sqrt2}$，正向
法向量 $N = \frac{(-1, 1)}{\sqrt2}$，而
\omterm{def:b2:curves:curvature}{曲率中心}为
\[
(0, 1) + 2\sqrt2\cdot\frac{(-1, 1)}{\sqrt2} = (-2,\ 3) :
\]
\omterm{def:b2:curves:curvature}{密切圆}的方程为 $(x + 2)^2 + (y - 3)^2 =
8$。作为对“最佳圆近似”这一说法的检验：
由圆的方程关于 $y$ 求解，并在 $(0,1)$ 附近展开，
得 $y = 1 + x + \frac{x^2}2 + O(x^3)$——
恰是 $\eu^x$ 的二阶泰勒展开式。
\omterm{def:b2:curves:curvature}{密切圆}与曲线的函数值、斜率\emph{以及}
二阶导数都相符；普通的切圆只能
与前两项相符。
\end{example}

\begin{example}[圆的渐屈线即
其圆心]\label{ex:b2:curves:circleevolute}
对按逆时针方向绕行、半径为 $R$ 的圆，
则 $\kappa = 1/R$，$N$ 指向圆心，故
\omterm{def:b2:curves:curvature}{曲率中心} $\tilde\gamma + \frac1\kappa N$ 即
圆心，对每个 $s$ 均成立：圆的\omterm{def:b2:curves:curvature}{密切圆}
即圆本身，\omterm{def:b2:curves:curvature}{曲率中心}的轨迹
退缩为一点。这一退化情形
可用来校验 \cref{exo:b2:curves:6}：在该题中渐屈线的
速度为 $-\frac{\kappa'}{\kappa^2}N$，它
恒为零，当且仅当 $\kappa$ 为常数。
\end{example}

\begin{proposition}[任意参数化
下的曲率]\label{prop:b2:curves:kappaformula}
对正则的 $\mathcal{C}^2$ 平面弧 $\gamma(t) = (x(t), y(t))$，
\[
\kappa(t)
= \frac{x'(t)\,y''(t) - y'(t)\,x''(t)}
       {\bigl(x'(t)^2 + y'(t)^2\bigr)^{3/2}} ,
\]
特别地，$\kappa = \dfrac{y''}{(1 + y'^2)^{3/2}}$ 适用于
图像 $y = f(x)$。
\end{proposition}

\begin{proof}
记 $v(t) = \norm{\gamma'(t)} = s'(t)$，于是 $\gamma' = vT$
（由单位速率情形与 $s$ 复合得到），求导得
\[
\gamma'' = v'T + v\,T'\cdot s' = v'T + v^2\kappa N .
\]
现在对标准有向基下的
$(\gamma', \gamma'')$ 取\omterm{def:b2:linalg:det}{行列式}：由于 $\det(T, T) = 0$，且 $\det(T, N) = 1$，
\[
\det(\gamma', \gamma'') = \det(vT,\ v'T + v^2\kappa N)
= v^3\kappa .
\]
左边为 $x'y'' - y'x''$，又 $v^3 = (x'^2 +
y'^2)^{3/2}$。图像情形即参数化 $t \mapsto (t,
f(t))$。
\end{proof}

\begin{example}[圆、直线、抛物线]\label{ex:b2:curves:kappaexamples}
直线满足 $\kappa = 0$（反之亦然：$T' = 0$ 表示 $T$
为常向量，故 $\tilde\gamma(s) = \tilde\gamma(0) + sT$，即一条直线）。
半径为 $R$、按逆时针方向绕行的圆有 $\kappa = 1/R$：
取 $\gamma(t) = (R\cos t, R\sin t)$，由公式得 $\kappa =
R^2/R^3$。对抛物线 $y = x^2/2$：$\kappa(x) = 1/(1 +
x^2)^{3/2}$，在顶点处取最大值——抛物线
在拐弯处弯得最厉害。
\end{example}

\begin{remark}[曲率的常见误区]
(i) 平面弧的代数\omterm{thm:b2:curves:frenet2d}{曲率}在弧或平面的定向被反转时
会变号；只有 $\abs\kappa$ 与 $R = 1/\abs\kappa$
是纯几何量。一个顺时针绕行的圆
有 $\kappa = -1/R$。(ii) 函数图像的公式
$\kappa = f''/(1 + f'^2)^{3/2}$ 悄悄取了以 $x$ 为
参数化方式；把它用于该点附近并非图像的曲线
（如竖直\omterm{def:b2:curves:arc}{切线}处）是经典的错误。
(iii) 在 $\gamma' = 0$ 的点处，
$T$ 与 $\kappa$ 都无定义，而
轨迹可能真正发生断裂
（见\cref{ex:b2:curves:cusp}）；求导前务必先检验
正则性，再对单位切向量求导。(iv) 在空间中，按约定 $\kappa =
\norm{T'} \geq 0$，因此既没有可能出错的符号，
也没有可利用的符号——拐点类信息
转入\omterm{thm:b2:curves:frenet3d}{挠率}。(v) 最后，$\kappa$ 是一个
\emph{对弧长的导数}：对非单位速率的参数化而言，
忘记 $v^3$ 这一因子，
在\cref{prop:b2:curves:kappaformula}中，
这是实践中最常见的错误。
\end{remark}

\begin{figure}[ht]
\centering
\begin{tikzpicture}[scale=1.6]
  % parabola y = x^2/2
  \draw[->] (-2.2,0) -- (2.4,0) node[below] {$x$};
  \draw[->] (0,-0.4) -- (0,2.4) node[left] {$y$};
  \draw[thick, ocDarkBlue, domain=-2.1:2.1, smooth, samples=60]
    plot (\x, {0.5*\x*\x});
  % osculating circle at vertex: kappa = 1 => R = 1, center (0,1)
  \draw[ocRed, dashed] (0,1) circle (1);
  \fill[ocRed] (0,1) circle (0.03) node[left, font=\footnotesize]
    {曲率中心};
  \fill (0,0) circle (0.03);
  % tangent and normal at vertex
  \draw[->, thick] (0,0) -- (0.7,0) node[below, font=\small] {$T$};
  \draw[->, thick] (0,0) -- (0,0.7) node[left, font=\small] {$N$};
  % a second point
  \fill (1.2,0.72) circle (0.03);
  \draw[->, thick] (1.2,0.72) -- ++(0.448,0.5376)
    node[right, font=\small] {$T$};
  \draw[->, thick] (1.2,0.72) -- ++(-0.5376,0.448)
    node[above, font=\small] {$N$};
\end{tikzpicture}
\caption{抛物线 $y = x^2/2$，其活动标架 $(T,
N)$ 及顶点处的\omterm{def:b2:curves:curvature}{密切圆}（半径 $1$，因为
$\kappa(0) = 1$）。标架随点的运动而转动；\omterm{thm:b2:curves:frenet2d}{曲率}
是单位弧长上这种转动的速率。}
\label{fig:b2:curves:osculating}
\end{figure}

\begin{theorem}[曲线由曲率决定]\label{thm:b2:curves:fundamental}
设 $\kappa \colon J \to \R$ 为\omterm{def:b2:metric:continuity}{连续}函数。则在平面上存在一条
单位速率的 $\mathcal{C}^2$ 弧，其\omterm{thm:b2:curves:frenet2d}{曲率}为
$\kappa$，且在相差一个直保距变换（先旋转
再平移）的意义下唯一。
\end{theorem}

\begin{proof}
\emph{存在性。}取 $s_0 \in J$，并令 $\varphi(s) =
\int_{s_0}^s \kappa(u)\,\dd u$，则
\[
\tilde\gamma(s) = \Bigl(\int_{s_0}^s \cos\varphi(u)\,\dd u,\
\int_{s_0}^s \sin\varphi(u)\,\dd u\Bigr).
\]
于是 $T(s) = (\cos\varphi(s), \sin\varphi(s))$ 是单位向量，
$N(s) = (-\sin\varphi, \cos\varphi)$，且
\[
T'(s) = \varphi'(s)\,(-\sin\varphi, \cos\varphi)
= \kappa(s)\,N(s) :
\]
因此该弧是单位速率的，\omterm{thm:b2:curves:frenet2d}{曲率}为 $\kappa$。

\emph{唯一性。}设 $\gamma_1, \gamma_2$ 为具有
相同\omterm{thm:b2:curves:frenet2d}{曲率}的单位速率弧。每个单位切向量都可以
借助一个显式构造提升为角函数：将 $T_j$ 视为复数
$z_j = a_j + \iu b_j$，其模为 $1$，取
$\varphi_j(0)$ 使 $z_j(0) = \eu^{\iu\varphi_j(0)}$，并令
\[
\varphi_j(s) = \varphi_j(0) + \int_0^s\det\bigl(T_j,
T_j'\bigr)(u)\,\dd u .
\]
由 $\abs{z_j} = 1$ 得 $\operatorname{Re}(\conj{z_j}\,z_j') =
0$，故 $\conj{z_j}\,z_j' = \iu\det(T_j, T_j') =
\iu\,\varphi_j'$，即 $z_j' = \iu\varphi_j'z_j$；于是
\[
\bigl(z_j\,\eu^{-\iu\varphi_j}\bigr)'
= \eu^{-\iu\varphi_j}\bigl(z_j' - \iu\varphi_j'z_j\bigr) = 0,
\]
于是处处有 $z_j = \eu^{\iu\varphi_j}$，即 $T_j =
(\cos\varphi_j, \sin\varphi_j)$，其中 $\varphi_j$ 属于
$\mathcal C^1$ 类。此外 $\det(T_j, T_j') = \det(T_j,
\kappa N_j) = \kappa$，故 $\varphi_j' = \kappa$。因此 $\varphi_2 = \varphi_1 + c$，
其中 $c$ 为常数：$T_2$ 是 $T_1$ 绕固定角 $c$ 旋转所得，
积分可得 $\gamma_2 = \rho(\gamma_1) + w$，其中 $\rho$ 为
旋转角为 $c$ 的旋转，$w$ 为常向量。
\end{proof}

\begin{remark}
这是\emph{几何基本定理}的一维原型：一组\omterm{def:b2:metric:complete}{完备}的
局部不变量（此处为一个函数）
在刚体运动下将对象分类。
下面的三维版本需要两个不变量。
\end{remark}

\begin{example}[常曲率
即圆]\label{ex:b2:curves:constantkappa}
在存在性的构造中取 $\kappa \equiv \kappa_0 > 0$：
$\varphi(s) = \kappa_0 s$，则
\[
\tilde\gamma(s) = \Bigl(\frac{\sin\kappa_0s}{\kappa_0},\
\frac{1 - \cos\kappa_0s}{\kappa_0}\Bigr) :
\]
这是半径为 $1/\kappa_0$、中心在 $(0,
1/\kappa_0)$ 的圆，按单位速率绕行。
由定理的唯一性部分，\emph{任一}
常\omterm{thm:b2:curves:frenet2d}{曲率} $\kappa_0$ 的单位速率弧都是半径为
$1/\kappa_0$ 的圆的一段（若 $\kappa_0 = 0$，则为直线）——
这正是\cref{ex:b2:curves:kappaexamples}中计算的逆，
以及\cref{exo:b2:curves:9}的平面情形。
\end{example}

\begin{example}[由曲率重建
曲线]\label{ex:b2:curves:reconstruction}
哪一条单位速率曲线的\omterm{def:b2:curves:curvature}{曲率半径}为 $R(s) = 1 +
s^2$？沿用存在性证明，取 $\kappa(s) =
\frac1{1+s^2}$，$s_0 = 0$：$\varphi(s) = \arctan s$，于是
\[
T(s) = (\cos\arctan s,\ \sin\arctan s)
= \Bigl(\frac{1}{\sqrt{1+s^2}},\
\frac{s}{\sqrt{1+s^2}}\Bigr),
\]
积分得
\[
\tilde\gamma(s) = \Bigl(\ln\bigl(s + \sqrt{1 + s^2}\bigr),\
\sqrt{1 + s^2} - 1\Bigr).
\]
令 $x = \ln(s + \sqrt{1+s^2})$，即 $s = \sinh x$，则
第二个坐标为 $\cosh x - 1$：该曲线是
\emph{悬链线} $y = \cosh x - 1$。这与
\cref{exo:b2:curves:3}首尾呼应，在那里我们直接算得 $R = \cosh^2 x = 1 +
\sinh^2 x = 1 + s^2$。由基本定理
保证悬链线是具有该\omterm{thm:b2:curves:frenet2d}{曲率}分布的\emph{唯一}曲线，
在相差一个保向等距的意义下。
\end{example}

\begin{remark}[曲率的后续应用]
分解 $\gamma'' = v'T + v^2\kappa N$ 是在
\cref{prop:b2:curves:kappaformula}的证明中得到的，它是每一种曲线运动的
运动学：切向
与向心加速度。\omterm{thm:b2:curves:frenet2d}{曲率}在曲面中再次出现，
（见\cref{ch:b2:surfaces}），通过画在其上的曲线的\omterm{thm:b2:curves:frenet2d}{曲率}
来体现；而本章周末问题中的\omterm{pb:b2:curves:1}{包络}理论——渐屈线与焦散线——
是波前的几何光学。
本科三年级分册将再次从内蕴观点出发
研究 $\R^n$ 中的子流形。
\end{remark}

\section{空间中的弗勒内标架}

现设 $\tilde\gamma \colon J \to \R^3$ 为单位速率的
$\mathcal{C}^3$ 弧，且它是\emph{双正则}的：$T'(s) \neq 0$
对所有 $s$ 成立。则 $\kappa(s) = \norm{T'(s)} > 0$ 定义
\emph{\omterm{thm:b2:curves:frenet2d}{曲率}}（空间中不带符号：不存在优选的法向
取向），并令：
\[
N(s) = \frac{T'(s)}{\kappa(s)}
\quad\text{（主法向量）},
\qquad
B(s) = T(s) \wedge N(s)
\quad\text{（副法向量）} ,
\]
于是 $(T, N, B)$ 为正向标准正交标架，即
\emph{弗勒内标架}。过 $\tilde\gamma(s)$ 且由
$T, N$ 张成的平面称为\emph{密切平面}。

\begin{theorem}[空间中的弗勒内公式]\label{thm:b2:curves:frenet3d}
存在\omterm{def:b2:metric:continuity}{连续}函数 $\tau$，称为\emph{挠率}\index{挠率}，使得
\[
T' = \kappa N, \qquad
N' = -\kappa T + \tau B, \qquad
B' = -\tau N .
\]
\end{theorem}

\begin{proof}
第一个公式即 $N$ 的定义。向量
$T, N, B$ 的模均为 $1$，且两两正交；
对六个关系式 $\langle X, Y\rangle = \delta_{XY}$ 求导，
可知 $(T', N', B')$ 在基 $(T, N, B)$ 下的矩阵是
反对称的：事实上 $\langle X', Y\rangle + \langle X, Y'\rangle = 0$，且
$\langle X', X\rangle = 0$。其 $(N, T)$ 元为 $\langle N', T\rangle =
-\langle N, T'\rangle = -\kappa$，其 $(T, B)$ 列元为
$\langle T', B\rangle = \kappa\langle N, B\rangle = 0$。记剩下的
自由元为 $\tau = \langle N', B\rangle$，即得上面给出的三个
公式：由反对称性推出 $\langle B', N\rangle = -\tau$ 与
$\langle B', T\rangle = 0$。$\tau = \langle N', B\rangle$ 的\omterm{def:b2:metric:continuity}{连续}性显然，
因为 $N'$ 与 $B$ \omterm{def:b2:metric:continuity}{连续}（$\tilde\gamma$ 为
$\mathcal{C}^3$，故 $N = T'/\kappa$ 为 $\mathcal{C}^1$）。
\end{proof}

\begin{example}[达布向量]\label{ex:b2:curves:darboux}
三个弗勒内公式可以合并为一个。令 $\omega(s) =
\tau\,T + \kappa\,B$（称为\emph{达布向量}）。利用 $B
\wedge T = N$，$T \wedge N = B$，$N \wedge B = T$：
\[
\omega \wedge T = \kappa\,N = T', \qquad
\omega \wedge N = \tau\,B - \kappa\,T = N', \qquad
\omega \wedge B = -\tau\,N = B' :
\]
每个标架向量都按 $X' = \omega \wedge X$ 演化，这正是
\emph{瞬时旋转}的运动学特征，其角速度向量为
$\omega$。标架以速率
$\norm\omega = \sqrt{\kappa^2 + \tau^2}$ 绕
动轴 $\omega$ 旋转；\omterm{thm:b2:curves:frenet2d}{曲率}是绕副法向量方向的分量，
\omterm{thm:b2:curves:frenet3d}{挠率}是绕切向方向的分量。对于
螺旋线，$\omega$ 是沿圆柱轴线方向的常向量——
这正是螺旋线的标架
匀速进动的原因。弗勒内矩阵的反对称性，在
\cref{exo:b2:curves:7}中被利用，它正是这一单一
几何事实的矩阵形式。
\end{example}

\begin{proposition}[挠率刻画平面性]\label{prop:b2:curves:torsion}
双正则弧包含于一个平面，当且仅当 $\tau
\equiv 0$；此时该平面即为（常值的）
密切平面。
\end{proposition}

\begin{proof}
若 $\tau \equiv 0$，则 $B' = 0$，故 $B$ 为常单位
向量 $B_0$，且
\[
\frac{\dd}{\dd s}\langle \tilde\gamma(s), B_0\rangle
= \langle T(s), B_0\rangle = 0 :
\]
$\langle \tilde\gamma, B_0\rangle$ 为常数，故弧位于一个平面内，
该平面与 $B_0$ 正交。反之，若弧位于平面 $P$ 内，
则 $T$ 与 $T'$（从而 $N$）都平行于
$P$ 的方向，对所有 $s$ 成立；故 $B = T \wedge N$ 为
$P$ 的两个单位法向量之一，由于它\omterm{def:b2:metric:continuity}{连续}，它为常向量；于是 $0 = B'
= -\tau N$，而 $N \neq 0$ 迫使 $\tau \equiv 0$。
\end{proof}

\begin{example}[倾斜圆的挠率为零]
弧 $\gamma(t) = \bigl(\cos t,\ \tfrac{\sin t}{\sqrt2},\
\tfrac{\sin t}{\sqrt2}\bigr)$ 位于平面 $y = z$ 内，且
是该平面内的单位圆（验证：$\norm{\gamma(t)} =
1$，该平面的标准正交基 $(1,0,0)$、
$(0, \tfrac1{\sqrt2}, \tfrac1{\sqrt2})$ 给出标准
参数化）。不作任何弗勒内计算，
由\cref{prop:b2:curves:torsion}可知 $\tau \equiv 0$，且
固定的副法向量必为该平面的单位法向量 $\pm(0,
\tfrac1{\sqrt2}, -\tfrac1{\sqrt2})$。\omterm{thm:b2:curves:frenet3d}{挠率}并不度量
“在空间中倾斜”，而是度量曲线\emph{离开}
某平面的程度。只有下面螺旋线中非零的 $b$ 才
产生真正的\omterm{thm:b2:curves:frenet3d}{挠率}。
\end{example}

\begin{example}[螺旋线]\label{ex:b2:curves:helixfrenet}
对螺旋线 $\gamma(t) = (a\cos t, a\sin t, bt)$，$a > 0$，
已算得 $s = ct$，其中 $c = \sqrt{a^2 + b^2}$。于是
\[
T = \frac1c(-a\sin t,\ a\cos t,\ b),
\qquad
T' \cdot \frac{\dd t}{\dd s}
= \frac{1}{c^2}(-a\cos t, -a\sin t, 0),
\]
因此 $\kappa = a/c^2 = a/(a^2 + b^2)$，且 $N = (-\cos t, -\sin t,
0)$：主法向量水平指向轴线。
接着 $B = T \wedge N = \frac1c(b\sin t, -b\cos t, a)$，且 $B'
\frac{\dd t}{\dd s} = \frac{b}{c^2}(\cos t, \sin t, 0) = -\tau N$
由此得
\[
\boxed{\ \kappa = \frac{a}{a^2 + b^2}, \qquad
\tau = \frac{b}{a^2 + b^2}. \ }
\]
两个不变量均为常数——并且可以证明，反之，
满足常数 $\kappa > 0$ 与常数 $\tau$ 的唯一双正则曲线
为螺旋线（$\tau = 0$ 时为圆）。注意符号：$b >
0$ 给出右手螺旋线，其\omterm{thm:b2:curves:frenet3d}{挠率}为正。
\end{example}

\begin{example}[不借助弧长参数的曲率与挠率；
扭立方曲线]\label{ex:b2:curves:twistedcubic}
通常无法以封闭形式按弧长重新参数化，
因此须由原始
导数直接求出不变量。记 $v = \norm{\gamma'} = s'$；则 $\gamma' =
vT$，并如
\cref{prop:b2:curves:kappaformula}的证明中所示，
\[
\gamma'' = v'T + v^2\kappa N,
\qquad
\gamma' \wedge \gamma'' = v^3\kappa\,(T \wedge N) =
v^3\kappa\,B .
\]
取\omterm{def:b2:nvs:norm}{范数}（空间中 $\kappa \geq 0$）：
\[
\kappa = \frac{\norm{\gamma' \wedge \gamma''}}{v^3} .
\]
再对 $\gamma''$ 求导，并将 $N' =
v(-\kappa T + \tau B)$ 代入（经 $s$ 的链式法则），唯一的
$B$ 分量仅来自最后一项：
\[
\gamma''' = \bigl(v'' - v^3\kappa^2\bigr)T +
\bigl(v'v\kappa + (v^2\kappa)'\bigr)N + v^3\kappa\tau\,B,
\]
于是，与 $\gamma' \wedge \gamma'' = v^3\kappa B$ 作内积，
\[
\det(\gamma', \gamma'', \gamma''') = \langle\gamma' \wedge
\gamma'',\ \gamma'''\rangle = v^6\kappa^2\tau,
\qquad\text{即}\qquad
\tau = \frac{\det(\gamma', \gamma'',
\gamma''')}{\norm{\gamma' \wedge \gamma''}^2} .
\]
将其应用于\emph{扭立方曲线} $\gamma(t) = (t,\
t^2,\ t^3)$ 在 $t = 0$ 处：$\gamma' = (1, 0, 0)$，$\gamma'' =
(0, 2, 0)$，$\gamma''' = (0, 0, 6)$，故 $v = 1$，
\[
\gamma' \wedge \gamma'' = (0, 0, 2), \qquad \kappa(0) = 2,
\qquad
\det(\gamma', \gamma'', \gamma''') = 12, \qquad \tau(0) =
\frac{12}{4} = 3 .
\]
小结：两个公式都是比值，其中速率
$v$ 恰以所需的次数约去——$\kappa$
像单位长度上的二阶导数那样缩放，$\tau$ 则
像三个导数的混合体积（每单位平方面积）那样缩放——这
正是它们具有几何意义的\emph{原因}，而 $\gamma''$ 本身
则不然。
\end{example}

\begin{remark}[空间曲线基本定理]
与平面情形类似，数对 $(\kappa, \tau)$，其中 $\kappa > 0$，
便在 $\R^3$ 的保向等距之下唯一确定一条双正则弧：
弗勒内公式是标架 $(T, N, B)$ 满足的线性\omterm{def:b2:diffcalc:differential}{微分}方程组，
可应用柯西–\omterm{def:b2:metric:continuity}{利普希茨}理论（参见
\cref{ch:b2:diffeq}）；解标架的正交性得以保持，
这是因为系数矩阵反对称（论证与格拉姆矩阵的情形相同，
见\cref{exo:b2:curves:7}），
曲线则由对 $T$ 积分恢复。细节留给
读者，作为一道内容较多而富有启发性的练习。
\end{remark}

\section{局部研究：相对于切线的位置}

\begin{proposition}[正则点处的
局部形状]\label{prop:b2:curves:local}
设 $\gamma$ 为 $\mathcal{C}^k$ 类的平面弧，
在 $t_0$ 处，$p$ 为使 $\gamma^{(p)}(t_0) \neq 0$ 的最小指标，$q$ 为
$> p$ 的最小指标，且使 $\gamma^{(q)}(t_0)$ 与
$\gamma^{(p)}(t_0)$ 不共线（假定二者均存在，$q \leq k$）。
在基 $(u, v) = (\gamma^{(p)}(t_0), \gamma^{(q)}(t_0))$
以 $\gamma(t_0)$ 为中心，由泰勒–杨公式得坐标
\[
X(t) \sim \frac{(t - t_0)^p}{p!},
\qquad
Y(t) \sim \frac{(t - t_0)^q}{q!} .
\]
局部图像仅取决于 $p$ 与 $q$ 的奇偶性：
\begin{center}
\begin{tabular}{c c l}
$p$ 奇，$q$ 偶 & \emph{普通点} & 曲线不穿过任何直线，
始终位于\omterm{def:b2:curves:arc}{切线}的一侧 \\
$p$ 奇，$q$ 奇 & \emph{拐点} & 曲线穿过其
\omterm{def:b2:curves:arc}{切线} \\
$p$ 偶，$q$ 奇 & \emph{第一类尖点} & 两支
分别位于\omterm{def:b2:curves:arc}{切线}的两侧 \\
$p$ 偶，$q$ 偶 & \emph{第二类尖点} & 两支
位于\omterm{def:b2:curves:arc}{切线}的同一侧 \\
\end{tabular}
\end{center}
\end{proposition}

\begin{proof}
用 $q$ 阶泰勒–杨公式（函数 $\gamma$ 为
$\mathcal{C}^q$ 类，在 $t_0$ 附近）：
\[
\gamma(t) - \gamma(t_0)
= \sum_{j=p}^{q} \frac{(t-t_0)^j}{j!}\,\gamma^{(j)}(t_0)
+ o\bigl((t-t_0)^q\bigr).
\]
由 $p$ 与 $q$ 的选取，每个 $\gamma^{(j)}(t_0)$（$p
\leq j < q$）都与 $u$ 共线；在
基 $(u, v)$ 下取分量，得 $X(t) = \frac{(t-t_0)^p}{p!}(1 + o(1))$ 且
$Y(t) = \frac{(t-t_0)^q}{q!}(1 + o(1))$。$X$ 与
$Y$ 的符号表（对 $t \gtrless t_0$）完全由奇偶性决定，
由此给出四种图形：例如当 $p$ 为偶数时，$X > 0$ 在
两侧都成立（两支都沿 $+u$ 方向离开，即尖点），
而\omterm{def:b2:curves:arc}{切线}一侧（$\operatorname{sign} Y$）随
$q$ 为奇数而翻转。
\end{proof}

\begin{example}
取 $\gamma(t) = (t^2, t^3)$，$t_0 = 0$
（见\cref{ex:b2:curves:cusp}）：$\gamma'' (0)= (2, 0)$，$\gamma'''(0) =
(0, 6)$，于是 $p = 2$，$q = 3$：第一类尖点，
即熟悉的半立方抛物线图形。取 $\gamma(t) = (t,
t^3)$ 在 $0$ 处：$p = 1$，$q = 3$：拐点——三次曲线穿过
其\omterm{def:b2:curves:arc}{切线}。
\end{example}

\begin{remark}[本册内的展望]
曲线以三种方式服务于后续各章。画在
曲面上时，它们确定其切平面与第一
基本形式（见\cref{ch:b2:surfaces}），其长度
通过限制环境度量来计算——
后一章基本上就是本章的相对化版本。
周末问题中的\omterm{pb:b2:curves:1}{包络}计算与二重
积分在\cref{ch:b2:multint}中相遇，\omterm{pb:b2:curves:1}{星形线}的面积
由格林公式重新算出
（见\cref{exo:b2:multint:5}）——同一条曲线，两种理论，
答案一致。此外，弗勒内方程组已经用到
了\cref{ch:b2:diffeq}中的线性\omterm{def:b2:diffcalc:differential}{微分}方程
（存在性、唯一性，以及保持正交性的论证，
见\cref{exo:b2:curves:7}）：曲线基本定理
本质上是一条\omterm{def:b2:diffcalc:differential}{微分}方程定理，
只是披着几何外衣。
\end{remark}

\begin{remark}[方法：局部分析的步骤]
实际操作中，分类是一套四步程序。
\emph{第一步}，在 $t_0$ 处求导，直到出现第一个非零
导数：其阶数为 $p$，其值为向量
$u$。\emph{第二步}，继续求导，直到出现一个
与 $u$ 不共线的导数：阶数为 $q$，向量为 $v$。
\emph{第三步}，在表中读取奇偶性 $(p, q)$。
\emph{第四步}，作图：曲线沿 $+u$ 方向离开（若 $p$ 为奇数；
否则先沿 $u$ 前进再沿 $u$ 折回，若 $p$ 为偶数），
位于 $v$ 的哪一侧由 $Y$ 的符号决定。两点需注意。
标架 $(u, v)$ 一般\emph{不}是标准正交的——表格
描述的是相对于\omterm{def:b2:curves:arc}{切线}的位置，而非角度
或距离，因此不能从图形中读出\omterm{thm:b2:curves:frenet2d}{曲率}。
此外，与 $u$ 共线的中间阶导数是允许的，
在阶 $p$ 与 $q$ 之间（它们只改变
$X$ 的展开）；必须\emph{避免}的做法是在第一个
非零导数处停下并猜测 $q = p + 1$：对于 $\gamma(t)
= (t^2, t^4 + t^5)$，这种朴素的猜测 $q = 3$ 是错误的，因为
$\gamma^{(3)}(0)$ 仍与 $\gamma''(0)$ 共线——
这正是\cref{exo:b2:curves:5}的内容。
\end{remark}

\section{练习}

\begin{exercise}[$\star$]\label{exo:b2:curves:1}
求摆线 $\gamma(t) = (t -
\sin t,\ 1 - \cos t)$，$t \in [0, 2\pi]$ 的一拱的弧长。\emph{（利用 $1 - \cos t =
2\sin^2(t/2)$。）}
\end{exercise}

\begin{exercise}[$\star$]\label{exo:b2:curves:2}
计算椭圆 $\gamma(t) = (a\cos t,\ b\sin
t)$（$a > b > 0$）的\omterm{thm:b2:curves:frenet2d}{曲率}，并求\omterm{thm:b2:curves:frenet2d}{曲率}最大与
最小的点。
\end{exercise}

\begin{exercise}[$\star$]\label{exo:b2:curves:3}
证明 $f(x) = \cosh x$ 的图像在
$[0, x]$ 上的弧长等于 $\sinh x$，并计算该
曲线（\emph{悬链线}）的\omterm{thm:b2:curves:frenet2d}{曲率}。验证 $R(x) = 1/\kappa(x) =
\cosh^2 x$。
\end{exercise}

\begin{exercise}[$\star\star$]\label{exo:b2:curves:4}
（对数螺线）设 $\gamma(t) = e^{t}(\cos t,\ \sin t)$，$t
\in \R$。
证明 $\gamma(t)$ 与 $\gamma'(t)$ 之间的夹角为
常数，计算 $\gamma$ 在 $(-\infty, 0]$ 上的弧长
（有限！），并计算\omterm{thm:b2:curves:frenet2d}{曲率}。
\end{exercise}

\begin{exercise}[$\star\star$]\label{exo:b2:curves:5}
求 $p$、$q$，并判断局部形状（普通点、拐点、
尖点）：$\gamma(t) = (t^2,\ t^4 + t^5)$ 在 $t = 0$ 处，以及
$\gamma(t) = (t^3,\ t^4)$ 在 $t = 0$ 处的情形。
\end{exercise}

\begin{exercise}[$\star\star$]\label{exo:b2:curves:6}
设 $\gamma$ 为单位速率的平面弧，且 $\kappa(s) > 0$ 对
所有 $s$ 成立，并令 $c(s) = \gamma(s) + \frac{1}{\kappa(s)}N(s)$ 为
\omterm{def:b2:curves:curvature}{曲率中心}（曲线 $c$ 称为\emph{渐屈线}）。
假设 $\kappa$ 为 $\mathcal{C}^1$ 类，证明 $c'(s) =
-\frac{\kappa'(s)}{\kappa(s)^2}N(s)$：渐屈线与
$\gamma$ 的法线相切。
\end{exercise}

\begin{exercise}[$\star\star\star$]\label{exo:b2:curves:7}
设 $A(s)$ 为一族\omterm{def:b2:metric:continuity}{连续}的反对称 $3 \times 3$
矩阵，$F' = F A$ 为一矩阵解，且 $F(s_0)$
为正交矩阵。证明 $F(s)$ 对所有 $s$ 均为正交矩阵。
\emph{（对 $G = F F^{\mathsf T}$ 求导，并利用柯西–\omterm{def:b2:metric:continuity}{利普希茨}
定理中的唯一性。）} 说明其与弗勒内方程组的关系。
\end{exercise}

\begin{exercise}[$\star\star\star$]\label{exo:b2:curves:8}
（闭凸曲线的总\omterm{thm:b2:curves:frenet2d}{曲率}）设 $\tilde\gamma$ 为
单位速率的 $\mathcal{C}^2$ 类闭平面弧，长为 $L$，
（即 $\tilde\gamma(s + L) = \tilde\gamma(s)$），
沿逆时针方向恰绕行一次。利用角度函数 $\varphi$（满足 $T =
(\cos\varphi, \sin\varphi)$，由
\cref{thm:b2:curves:fundamental}给出），说明为何 $\varphi(L) -
\varphi(0)$ 为 $2\pi$ 的整数倍，并证明
$\int_0^L \kappa(s)\,\dd s = \varphi(L) - \varphi(0)$。（对半径为
$R$ 的圆：$\int \kappa = \frac1R \cdot 2\pi R = 2\pi$。
\omterm{def:b2:curves:arc}{切线}转角定理断言，对每条简单闭曲线，其值为 $2\pi$，
本题不要求证明这一点。）
\end{exercise}

\begin{exercise}[$\star\star\star$]\label{exo:b2:curves:9}
证明：双正则空间曲线若其\omterm{thm:b2:curves:frenet2d}{曲率} $\kappa > 0$ 为常数，且
$\tau = 0$，则该曲线（或其一段）为半径为 $1/\kappa$ 的圆。
\emph{（利用\cref{prop:b2:curves:torsion}，再证明中心
$\gamma + \frac1\kappa N$ 保持不变。）}
\end{exercise}

\begin{exercise}[$\star$]\label{exo:b2:curves:10}
求抛物线 $y = x^2/2$ 在
$\intcc0a$ 上的弧长，并证明其等于
\[
\tfrac12\Bigl(a\sqrt{1 + a^2} + \ln\bigl(a + \sqrt{1 +
a^2}\bigr)\Bigr).
\]
\end{exercise}

\begin{exercise}[$\star\star$]\label{exo:b2:curves:11}
设 $\gamma$ 为正则的 $\mathcal C^2$ 类弧，它位于 $\R^n$ 中，
且其所有\omterm{def:b2:curves:arc}{切线}都经过定点 $P$。证明
$\gamma$ 的轨迹包含于某一条直线中。
\emph{（以弧长参数化，写出 $\gamma(s) +
\lambda(s)T(s) = P$ 并求导。）}
\end{exercise}

\begin{exercise}[$\star\star\star$]\label{exo:b2:curves:12}
（空间曲线基本定理）设 $\kappa > 0$ 与
$\tau$ 为区间 $J$ 上的\omterm{def:b2:metric:continuity}{连续}函数。完成
\cref{ex:b2:curves:helixfrenet}之后注记中所述的方案：
(a) 证明
线性方程组 $F' = FA(s)$，其中 $A(s)$ 为反对称的弗勒内
矩阵（由 $\kappa, \tau$ 构成），$F(s_0)$ 为直接
正交标架，该方程组具有唯一的整体解，且此解
始终为直接正交标架；(b) 构造单位速率的
双正则曲线，其\omterm{thm:b2:curves:frenet2d}{曲率}为 $\kappa$、\omterm{thm:b2:curves:frenet3d}{挠率}为 $\tau$；
(c) 证明在 $\R^3$ 的直接等距变换意义下唯一。
\end{exercise}

\section{周末问题：包络——星形线、两条渐屈线与
一条焦散线}

\begin{omfigure}
\begin{tikzpicture}[scale=2.2]
  \draw[->] (-1.15,0) -- (1.25,0) node[below] {$x$};
  \draw[->] (0,-1.15) -- (0,1.25) node[left] {$y$};
  \draw[omDef] (1,0) -- (0,0.0) (0.966,0) -- (0,0.259)
    (0.866,0) -- (0,0.5) (0.707,0) -- (0,0.707)
    (0.5,0) -- (0,0.866) (0.259,0) -- (0,0.966);
  \draw[omThm, very thick]
    plot[domain=0:360, samples=120]
    ({cos(\x)*cos(\x)*cos(\x)}, {sin(\x)*sin(\x)*sin(\x)});
\end{tikzpicture}

{\small 一架长为 $1$ 的梯子沿墙滑落时（蓝色
位置）从不穿过\omterm{pb:b2:curves:1}{星形线} $x^{2/3} + y^{2/3} = 1$，
（红色）：\omterm{pb:b2:curves:1}{星形线}是线段族的\emph{\omterm{pb:b2:curves:1}{包络}}，
与其中每一条线段都相切。}
\end{omfigure}

\begin{problem}[{周末问题——包络机与
四条经典曲线}]\label{pb:b2:curves:1}
单参数直线族通常无法均匀覆盖
平面：这些直线会沿着一条与它们全部相切的曲线堆积，
这条曲线即它们的\emph{包络}\index{包络}。光线使
\omterm{pb:b2:curves:1}{包络}以\emph{焦散线}的形式显现——咖啡杯里
那条明亮的带尖点曲线。本题先构造一般的\omterm{pb:b2:curves:1}{包络}机，
再用它依次处理四个例子：沿墙滑动的
梯子（\omterm{pb:b2:curves:1}{星形线}）、抛物线与摆线的法线，
（渐屈线，末尾附惠更斯摆），以及咖啡杯
中的焦散线（\omterm{pb:b2:curves:1}{肾形线}）。全文中，$D_t$ 表示
方程为 $a(t)\,x + b(t)\,y = c(t)$ 的直线，其中 $a, b, c$
为 $\mathcal C^2$ 类函数，且 $(a(t), b(t)) \neq (0,0)$，
并记 $\Delta(t) = a(t)b'(t) - a'(t)b(t)$。

\medskip
\noindent\textbf{第一部分——\omterm{pb:b2:curves:1}{包络}机。}
\begin{enumerate}
  \item 设 $\Delta(t) \neq 0$，证明
        \emph{特征方程组}
        \[
        \begin{cases} a(t)\,x + b(t)\,y = c(t)\\
        a'(t)\,x + b'(t)\,y = c'(t)\end{cases}
        \]
        有唯一解 $E(t) = (x(t), y(t))$，由下式给出：
        $x = \dfrac{cb' - c'b}{\Delta}$, $y = \dfrac{ac' -
        a'c}{\Delta}$.
  \item 进一步假设 $E$ 为 $\mathcal C^1$ 类（在 $t$ 附近），
        且 $E'(t) \neq 0$。对方程组第一式求导，
        并证明 $a(t)\,x'(t) +
        b(t)\,y'(t) = 0$，并由此推出曲线 $E$
        经过 $D_t$ 上的一点，且\emph{与
        $D_t$ 方向相同}：该直线族与 $E$ 相切，
        称其为该直线族的\emph{\omterm{pb:b2:curves:1}{包络}}。
  \item 验算：抛物线 $y =
        x^2/2$ 在点 $(t, t^2/2)$ 处的\omterm{def:b2:curves:arc}{切线}为 $tx - y =
        t^2/2$。验证\omterm{pb:b2:curves:1}{包络}机求出的正是
        抛物线本身。
  \item （法线族的\omterm{pb:b2:curves:1}{包络}）设 $\gamma$ 为单位速率曲线，
        且 $\kappa(s) \neq 0$。点 $\gamma(s)$ 处的法线
        为 $\{M : \langle M - \gamma(s), T(s)
        \rangle = 0\}$。证明其特征方程组
        蕴含 $\langle M - \gamma(s), N(s)\rangle =
        1/\kappa(s)$，从而特征点为
        \omterm{def:b2:curves:curvature}{曲率中心}：\emph{法线族的
        \omterm{pb:b2:curves:1}{包络}即为渐屈线}，从而恢复
        \cref{exo:b2:curves:6}。检验 $\Delta(s) =
        \kappa(s)$。
  \item 两种退化情形。对于直线束 $D_\theta : x\cos
        \theta + y\sin\theta = 0$，证明
        对每个 $\theta$，特征点均为原点
        （“\omterm{pb:b2:curves:1}{包络}”退化为一点，且 $E' = 0$：
        第 2 小题不适用）。对平行直线族
        （$a, b$ 为常数），证明 $\Delta \equiv 0$，且
        特征方程组一般不相容：
        即不存在\omterm{pb:b2:curves:1}{包络}。
\end{enumerate}

\medskip
\noindent\textbf{第二部分——滑动的梯子与
\omterm{pb:b2:curves:1}{星形线}。} 长为 $1$ 的线段滑动，其一端 $P_t =
(\cos t, 0)$ 在地面上，另一端 $Q_t = (0, \sin t)$ 在墙上，
$t \in \intoo0{\pi/2}$。
\begin{enumerate}[resume]
  \item 证明直线 $(P_tQ_t)$ 的方程为 $x\sin t +
        y\cos t = \sin t\cos t$，并由\omterm{pb:b2:curves:1}{包络}求法
        得出特征点
        \[
        E(t) = (\cos^3 t,\ \sin^3 t) :
        \]
        即\emph{星形线}\index{星形线}，其隐式方程为
        $x^{2/3} + y^{2/3} = 1$（通过对称延拓至
        其余象限）。
  \item 证明 $E'(0) = 0$，并利用局部分类
        （\cref{prop:b2:curves:local}），断定\omterm{pb:b2:curves:1}{星形线}
        在 $(1, 0)$ 处有一个第一类尖点
——其四个轴点处亦然。
  \item 计算 $\norm{E'(t)} = \tfrac32\abs{\sin 2t}$，并
        推出\omterm{pb:b2:curves:1}{星形线}的总长为 $6$。
  \item 梯子在何处与\omterm{pb:b2:curves:1}{星形线}相切？证明 $E(t) =
        P_t + \sin^2 t\,(Q_t - P_t)$，即切点
        将梯子按比 $\sin^2 t : \cos^2
        t$ 分成两段，
        随梯子滑动，切点由梯子一端扫向另一端。
  \item 计算\omterm{pb:b2:curves:1}{星形线}所围面积，证明
        第一象限面积为 $3\int_0^{\pi/2}\sin^4
        t\cos^2 t\,\dd t$，再用降幂法计算此积分
        （$\sin^2 2t = \tfrac{1 - \cos 4t}2$），
        并得出总面积为 $3\pi/8$。
\end{enumerate}

\medskip
\noindent\textbf{第三部分——抛物线的渐屈线。}
设 $\gamma(t) = (t, t^2/2)$。
\begin{enumerate}[resume]
  \item 证明 $\gamma(t)$ 处的法线方程为
        $x + t\,y = t + t^3/2$。
  \item 运用\omterm{pb:b2:curves:1}{包络}求法，证明法线族的
        \omterm{pb:b2:curves:1}{包络}为
        \[
        E(t) = \Bigl(-t^3,\ 1 + \tfrac32 t^2\Bigr),
        \]
        隐式方程为 $x^2 = \tfrac8{27}(y - 1)^3$，
        即半立方抛物线。
  \item 与第 4 题互相印证：计算
        \omterm{def:b2:curves:curvature}{曲率中心} $\gamma(t) + \frac1{\kappa(t)}N(t)$，由
        $\kappa(t) = (1 + t^2)^{-3/2}$
        （见\cref{ex:b2:curves:kappaexamples}）算出，并得到
        同一点。
  \item 证明渐屈线在
        $(0, 1)$ 处有第一类尖点，它即顶点处的\omterm{def:b2:curves:curvature}{曲率中心}
——也就是 $\kappa$ 取极值之处，这与
        $c' = -\frac{\kappa'}{\kappa^2}N$ 的公式
        （见\cref{exo:b2:curves:6}）所预言的结果相符。
  \item 抛物线的法线中有多少条通过给定点
        $(x_0, y_0)$？证明答案由
        三次方程 $\tfrac{t^3}2 + (1 - y_0)\,t - x_0 = 0$ 决定；
        完整讨论轴情形 $x_0 = 0$（$y_0 < 1$ 时一条法线，
        $y_0 > 1$ 时三条），并
        将渐屈线解释为分界曲线。
\end{enumerate}

\medskip
\noindent\textbf{第四部分——咖啡杯焦散线。} 平行光线
方向为 $(1, 0)$，射向镜面圆 $x^2 + y^2 = 1$ 的内侧；
射到 $P_\theta =
(\cos\theta, \sin\theta)$ 的光线按反射定律
反射。
\begin{enumerate}[resume]
  \item 由镜面关于法线（即半径）的对称性说明反射方向
        为 $v = u -
        2\langle u, n\rangle n$，其中 $u = (1,0)$，$n =
        (\cos\theta, \sin\theta)$，并计算 $v =
        -(\cos2\theta, \sin2\theta)$。
  \item 证明反射光线位于直线
        \[
        x\sin 2\theta - y\cos 2\theta = \sin\theta .
        \]
  \item 运用\omterm{pb:b2:curves:1}{包络}求法（$\Delta = 2$），证明
        焦散线为
        \[
        E(\theta) = \Bigl(\tfrac{3\cos\theta -
        \cos3\theta}4,\ \tfrac{3\sin\theta -
        \sin3\theta}4\Bigr),
        \]
        即\emph{肾形线}\index{肾形线}。
  \item 计算 $E'(\theta) = \tfrac32\sin\theta\,
        (\cos2\theta, \sin2\theta)$，验证\omterm{def:b2:curves:arc}{切线}
        方向即反射光线的方向（第 16 题），
        求两个尖点 $(\pm\tfrac12, 0)$，并
        证明反射光线与轴 $y = 0$ 相交于
        $x = \frac1{2\cos\theta}$ 处——因此近轴光线
        汇聚于 $x = \tfrac12$，即焦距为 $R/2$ 的
        镜面（半径为 $R$）。
  \item 证明\omterm{pb:b2:curves:1}{肾形线}总长为 $6$，并证明
        在 $\theta = 0$ 附近，
        \[
        E(\theta) - \bigl(\tfrac12, 0\bigr) =
        \bigl(\tfrac34\theta^2 + o(\theta^2),\ \theta^3 +
        o(\theta^3)\bigr) :
        \]
        即一个沿轴指向的第一类尖点。
  \item 用一段话说明焦散线为何明亮：
        焦散线外侧紧邻的每一点都有两条
        反射光线经过，而焦散线上每一点处的光线
        是“无穷集中”（映射 $(\theta,
        \text{沿光线的距离}) \mapsto \R^2$ 的
        临界点恰好位于\omterm{pb:b2:curves:1}{包络}上）。
\end{enumerate}

\medskip
\noindent\textbf{第五部分——惠更斯：摆线是它自身的
渐屈线。} 设 $\gamma(t) = (t - \sin t,\ 1 - \cos t)$，$t \in
\intoo0{2\pi}$，为摆线的一拱。
\begin{enumerate}[resume]
  \item 计算 $\kappa(t) = -\dfrac1{4\sin(t/2)}$ 及
        \omterm{def:b2:curves:curvature}{曲率中心}；证明渐屈线为
        \[
        c(t) = (t + \sin t,\ \cos t - 1),
        \]
        并证明代换 $t = u + \pi$ 表明它是
        原摆线经 $(\pi, -2)$ 平移所得的曲线：
        \emph{摆线的渐屈线是与之全等的
        摆线}（惠更斯）。
  \item 验证在顶点 $t =
        \pi$ 处的\omterm{def:b2:curves:curvature}{曲率半径}等于 $4$，恰为一拱长 $8$ 的一半
        （见\cref{exo:b2:curves:1}）；定出渐屈线的尖点，
        它位于顶点正下方，距离为 $4$。
  \item （弦性质）设 $\gamma$ 为单位速率曲线，且
        $\kappa > 0$，$\kappa$ 属于 $\mathcal C^1$ 类，并且
        $R = 1/\kappa$ 严格单调。利用 $c' = R'N$，
        证明渐屈线在
        $c(s_0)$ 与 $c(s_1)$ 之间的弧长为 $\abs{R(s_1) - R(s_0)}$。
        解释：将一根绷紧的细绳从渐屈线上展开，
        起始长度为 $R(s_0)$，其自由端
        描出原曲线——因此，在惠更斯摆钟的两片
        摆线形挡板之间摆动的摆，
        其轨迹为一条摆线。
  \item 综合。第一部分中的机器给出了
        \omterm{pb:b2:curves:1}{星形线}、半立方抛物线、\omterm{pb:b2:curves:1}{肾形线}和
        摆线。对这四族中的每一族，用一句话说明
        假设 $\Delta \neq 0$ 与 $E' \neq 0$ 在何处成立、何处失效，
        以及它们失效时所对应的几何事件
        （尖点、焦点、退化）；每次 $E' \neq 0$ 失效都预示着其中一种。
        \omterm{thm:b2:curves:frenet2d}{曲率}的极值点必定出现在
        法线\omterm{pb:b2:curves:1}{包络}的何处？并说明理由。
\end{enumerate}
\end{problem}
